\documentclass[11pt]{article}

\usepackage[T1]{fontenc}
\usepackage{lmodern}
\usepackage[a4paper,margin=28mm]{geometry}
\usepackage{amsmath,amssymb,amsthm,mathtools,mathrsfs}
\usepackage{microtype}
\usepackage{enumitem}
\usepackage{xcolor}
\usepackage[colorlinks=true,linkcolor=blue!45!black,citecolor=blue!45!black,
  urlcolor=blue!55!black]{hyperref}
\setlength{\emergencystretch}{2em}

\newtheorem{theorem}{Theorem}[section]
\newtheorem{proposition}[theorem]{Proposition}
\newtheorem{lemma}[theorem]{Lemma}
\newtheorem{corollary}[theorem]{Corollary}
\newtheorem*{maintheorem}{Main Theorem}
\theoremstyle{definition}
\newtheorem{definition}[theorem]{Definition}
\theoremstyle{remark}
\newtheorem{remark}[theorem]{Remark}

\newcommand{\C}{\mathbb C}
\newcommand{\N}{\mathbb N}
\newcommand{\cA}{\mathcal A}
\newcommand{\cD}{\mathcal D}
\newcommand{\cF}{\mathcal F}
\newcommand{\cK}{\mathcal K}
\newcommand{\cM}{\mathcal M}
\newcommand{\cP}{\mathcal P}
\newcommand{\ord}{\operatorname{ord}}
\newcommand{\supp}{\operatorname{supp}}
\newcommand{\Frac}{\operatorname{Frac}}
\newcommand{\trdeg}{\operatorname{trdeg}}

\title{A Positive Hypertranscendental Frame\\
for the Riemann Divisor}
\author{Michael Haufschild}
\date{}

\begin{document}

\maketitle

\begin{abstract}
Write \(M(n)=\max_p v_p(n)\) for \(n\geq2\), and let \(a_m\) be the indicator of
\(M(n)=m\).  The Dirichlet series
\[
 S_m(s)=\sum_{n\geq1}\frac{a_m(n)}{n^s}
\]
have nonnegative, pairwise disjoint coefficient supports and partition the zeta coefficients
away from \(n=1\).  We prove that this elementary partition has an unexpectedly rigid analytic
and differential structure.  For
\[
 H_m(s)=\frac1{\zeta((m+1)s)}-\frac1{\zeta(ms)}
\]
one has \(S_m=\zeta H_m\) on \(\Re s>1\).  If \(m\geq12\), then \(H_m\) is holomorphic and
zero-free throughout \(\Re s>1/2\).  Consequently every \(S_m\) carries exactly the Riemann
zeta zero divisor, including multiplicities, in that half-plane.

The common divisor is the only differential-algebraic coupling.  Let
\(\C(s)\langle\zeta\rangle\) be the differential field generated by the rational coordinate,
\(\zeta\), and all derivatives of \(\zeta\).  Independently of the zero-free threshold, we prove
the stronger algebraic law
\[
 \{H_m:m\geq2\}
\quad\text{is jointly differentially algebraically independent over}\quad
 \C(s)\langle\zeta\rangle .
\]
Restricting this larger independent family to \(m\geq12\) shows that a countable family of
positive, coefficient-disjoint Dirichlet series can have one complete common zero divisor while
its zero-free transverse factors have maximal possible differential rank even after that divisor
is adjoined.

The mechanism is a general \(\Omega\)-leading-echelon law for Dirichlet series.  Prime
derivations lower the least total-prime-factor grade of level \(m\) to \(m-1\); the lowest layer
is a single prime ray, and its finite-jet Jacobian is a product of prime-log Vandermonde
determinants.  Higher \(\Omega\)-layers cannot cancel this coefficient.  Applied simultaneously
to the constant zeta source and all exponent shells, the law yields the relative theorem.
It also gives an exact positive reformulation of the Riemann hypothesis: an off-critical zero
would have to occur, with the same multiplicity, in every one of infinitely many disjoint
arithmetic channels, although their normalized channels obey no finite algebraic differential
relation.  No hypothesis on zeta zeros is used.
\end{abstract}

\medskip
\noindent\textbf{2020 Mathematics Subject Classification.}
Primary 11M26, 11J91; Secondary 12H05, 11M41, 11A51.

\noindent\textbf{Keywords.}
Riemann zeta function; Riemann hypothesis; hypertranscendence; differential algebra;
Dirichlet series; \(k\)-free integers; arithmetic functions.

\section{Introduction}\label{sec:intro}

The zero set of a Dirichlet series is produced by cancellation among its coefficients.  For the
Riemann zeta function those coefficients are all equal to one, so the cancellation is global
and arithmetically unstructured.  This paper gives an exact decomposition into infinitely many
positive arithmetic channels for which the same zero divisor is reproduced separately in every
channel.

For \(n\geq2\), let
\[
 M(n)=\max_{p\mid n}v_p(n)
\]
be the largest exponent in the prime factorization of \(n\), and put \(M(1)=0\).  For \(m\geq1\)
define
\[
 a_m(n)=\mathbf 1_{\{M(n)=m\}},
 \qquad
 S_m(s)=\sum_{n\geq1}a_m(n)n^{-s}.
 \tag{1.1}
\]
Each coefficient belongs to exactly one shell:
\[
 a_m(n)\in\{0,1\},\qquad
 a_m(n)a_k(n)=0\quad(m\ne k),\qquad
 \sum_{m\geq1}a_m(n)=1-\mathbf 1_{\{1\}}(n).
 \tag{1.2}
\]
The familiar Dirichlet series of \(m\)-free integers then gives
\[
 S_m(s)=\zeta(s)H_m(s),\qquad
 H_m(s)=\zeta((m+1)s)^{-1}-\zeta(ms)^{-1}.
 \tag{1.3}
\]

Formula \eqref{eq:factor} below is elementary; by itself it says little about zeros because
\(H_m\) could introduce or remove them.  The first analytic result is that it does neither:
for every \(m\geq12\), \(H_m\) is holomorphic and zero-free on the full strict right half-plane
\(\Re s>1/2\).  Hence the positive subseries supported on each set
\(\{n:M(n)=m\}\) has exactly the zeta zero divisor there.  A hypothetical zero
\(\rho\) with \(1/2<\Re\rho<1\) is therefore not one cancellation of the all-ones series.  It is
simultaneously a cancellation in every one of the pairwise disjoint positive coefficient
streams \(a_m\), \(m\geq12\), with its analytic multiplicity unchanged.  Here and below
``cancellation'' means a zero of the meromorphic continuation of that positive-coefficient
Dirichlet series; no termwise convergence in the critical strip is asserted.

Common divisors usually signal algebraic dependence.  Here the opposite extreme occurs after
the common factor is removed.  Let
\[
 K_\zeta=\C(s)\langle\zeta\rangle
 =\C(s)\bigl(\zeta,\zeta',\zeta'',\ldots\bigr)
 \tag{1.4}
\]
inside a differential field of meromorphic germs on a right half-plane.  Our main theorem states
that the complete set
\[
 \{H_m^{(j)}:m\geq2,\ j\geq0\}
 \tag{1.5}
\]
is algebraically independent over \(K_\zeta\).  This is the strongest possible
differential-algebraic separation of the transverse factors: no polynomial differential
identity involving finitely many shells survives even when \(\zeta\), all of its derivatives,
and rational functions of \(s\) are allowed as coefficients.  We use the standard differential
field terminology of Kolchin~\cite{Kolchin1973}.

In fact the differential theorem has the larger natural range \(m\geq2\), although the
elementary zero-free estimate is asserted only for \(m\geq12\).  The result is not a consequence
of separate hypertranscendence.  Individual
hypertranscendence of broad classes of Dirichlet series belongs to a classical line beginning
with Ostrowski~\cite{Ostrowski1920}; modern extensions include general Dirichlet and automorphic
\(L\)-series~\cite{Nagoshi2016}.  Nor is equality of zero divisors alone new: Grosswald and
Schnitzer constructed modified zeta and \(L\)-functions preserving zeros and
multiplicities~\cite{GrosswaldSchnitzer1978}.  Jacobian methods for algebraic independence of
Dirichlet series were developed by Shapiro and Sparer~\cite{ShapiroSparer1986} and in
differential-algebraic form by Postnikov and Pong~\cite{Postnikov1949,Pong2016}.  The
contribution here is the simultaneous configuration:
\[
\begin{gathered}
\text{positive and pairwise disjoint coefficient supports},\\
\text{exact coefficient fusion to the zeta stream},\\
\text{one common right-half-plane divisor with multiplicity},\\
\text{countably maximal differential rank transverse to that divisor}.
\end{gathered}
\tag{1.6}
\]
The last property is relative to \(K_\zeta\), not merely to \(\C\), and is proved for the
explicit adjacent reciprocal-dilation differences in \eqref{eq:factor}.

We say that a family \((f_i)\) is \emph{jointly differentially algebraically independent} over a
differential field \(K\) when the ordinary family of all jets
\(\{f_i^{(j)}:i,\ j\geq0\}\) is algebraically independent over \(K\).

\begin{maintheorem}\label{thm:main-intro}
The exponent-shell family has the following unconditional structure.
\begin{enumerate}[label=\textup{(\alph*)}]
\item The \(a_m\) are \(0/1\)-valued with pairwise disjoint supports, and
\(\sum_{m\geq1}a_m=\mathbf1-\mathbf1_{\{1\}}\) coefficientwise.
\item For every \(m\geq1\), \(S_m=\zeta H_m\) initially on \(\Re s>1\).  For every
\(m\geq12\), \(H_m\) extends holomorphically and without zeros to \(\Re s>1/2\), so
\(S_m^\#\) and \(\zeta\) have identical zero divisors there, with multiplicity.
\item The larger family \((H_m)_{m\geq2}\) is jointly differentially algebraically independent
over \(\C(s)\langle\zeta\rangle\).
\item RH is equivalent to zero-freeness of one \(S_m^\#\), of every \(S_m^\#\), or of their
common zero set in \(1/2<\Re s<1\), where \(m\geq12\).
\end{enumerate}
\end{maintheorem}

The proof introduces a reusable source criterion.  Grade the Dirichlet-convolution algebra by
\[
 \Omega(n)=\sum_pv_p(n).
 \tag{1.7}
\]
A level-\(m\) shell begins in grade \(m\), and equality occurs only on pure powers \(p^m\).
After a prime derivation, the first layer is therefore a point mass at \(p^{m-1}\).
For any finite rectangle of logarithmic jets, distinct levels separate the determinant into
blocks and distinct prime logarithms make each block Vandermonde.  Since all omitted terms have
higher \(\Omega\)-grade, they cannot cancel the lowest determinant coefficient.  A Jacobian
criterion yields algebraic independence in the Dirichlet-convolution domain, and an
exponential-boundary uniqueness lemma transfers it over the genuine coordinate field
\(\C(s)\).

This construction gives progress toward RH in a precise structural sense, not a claim that RH is
proved.  It replaces one undifferentiated cancellation problem by a positive
\emph{Riemann-divisor frame}: infinitely many disjoint arithmetic channels carry the same
unknown divisor, while every finite differential-algebraic coupling among their normalized
directions is excluded.  Thus simultaneous shell estimates are genuinely independent analytic
constraints rather than disguised differential restatements of one channel.  The exact logical
RH consequence is stated in Corollary~\ref{cor:rh}; the theorem does not supply the additional
inequality needed to exclude their common zeros.

Section~\ref{sec:shells} proves the positive partition and factorization.
Section~\ref{sec:zerofree} proves the zero-free multiplier law by prime-two domination.
Section~\ref{sec:omega} establishes the universal \(\Omega\)-leading-echelon criterion.
Section~\ref{sec:application} applies it to zeta and the shells and proves relative
hypertranscendence.  Section~\ref{sec:rhframe} records the resulting divisor frame and its exact
RH content.

\section{Exponent shells and exact positive fusion}\label{sec:shells}

For \(r\geq1\), call \(n\) \(r\)-free if no \(p^r\) divides \(n\), and write
\[
 b_r(n)=\mathbf 1_{\{n\ {\rm is}\ r\text{-free}\}}.
\]
The convention makes \(b_1=\mathbf 1_{\{1\}}\).  Since \(r\)-freeness implies
\((r+1)\)-freeness,
\[
 a_m=b_{m+1}-b_m
 \label{eq:shell-difference}
 \tag{2.1}
\]
takes only the values zero and one.

\begin{proposition}[Positive shell partition]\label{prop:partition}
For every \(n\geq1\) and \(m,k\geq1\):
\begin{enumerate}[label=\textup{(\roman*)}]
\item \(a_m(n)=1\) if and only if \(M(n)=m\);
\item if \(m\ne k\), then \(a_m(n)a_k(n)=0\);
\item
\[
 \sum_{m\geq1}a_m(n)=1-\mathbf 1_{\{1\}}(n),
 \label{eq:pointwise-fusion}
 \tag{2.2}
\]
and the sum has finite support for each fixed \(n\);
\item for every \(R\geq1\),
\[
 \sum_{m=1}^{R}a_m(n)=b_{R+1}(n)-b_1(n).
 \label{eq:finite-telescope}
 \tag{2.3}
\]
\end{enumerate}
\end{proposition}

\begin{proof}
If \(n\geq2\), then \(n\) is \((m+1)\)-free but not \(m\)-free exactly when the largest
valuation of \(n\) is \(m\).  This proves (i) and (ii).  Formula \eqref{eq:finite-telescope}
is the telescoping sum of \eqref{eq:shell-difference}.  Taking
\(R\geq M(n)\) proves \eqref{eq:pointwise-fusion}; for \(n=1\), every difference is zero.
\end{proof}

The Dirichlet series of \(b_r\) is obtained directly from Euler factors:
\[
 \sum_{n\geq1}\frac{b_r(n)}{n^s}
 =\prod_p(1+p^{-s}+\cdots+p^{-(r-1)s})
 =\frac{\zeta(s)}{\zeta(rs)}
 \qquad(\Re s>1),
 \label{eq:mfree-series}
 \tag{2.4}
\]
where for \(r=1\) both sides equal one.

\begin{proposition}[Shell factorization]\label{prop:factor}
For every \(m\geq1\) and \(\Re s>1\),
\begin{equation}\label{eq:factor}
 S_m(s)=\zeta(s)H_m(s),\qquad
 H_m(s)=\frac1{\zeta((m+1)s)}-\frac1{\zeta(ms)}.
 \tag{2.5}
\end{equation}
Moreover the coefficient fusion in Proposition~\ref{prop:partition} is absolutely summable on
every half-plane \(\Re s\geq1+\varepsilon\).
\end{proposition}

\begin{proof}
Subtract \eqref{eq:mfree-series} at \(r=m+1\) and \(r=m\).  For
\(\sigma=\Re s>1\),
\[
 \sum_{m\geq1}\sum_{n\geq1}|a_m(n)n^{-s}|
 =\sum_{n\geq2}n^{-\sigma}<\infty,
\]
which proves the final assertion.
\end{proof}

In particular the countably infinite relation
\[
 \sum_{m\geq1}S_m(s)=\zeta(s)-1,\qquad \Re s>1,
 \label{eq:analytic-fusion}
 \tag{2.6}
\]
is exact.  The differential-independence theorem below concerns finite differential
polynomials, so \eqref{eq:analytic-fusion} is compatible with maximal finite algebraic rank:
an infinite locally convergent sum is not a differential-polynomial relation.

\section{A zero-free adjacent-dilation law}\label{sec:zerofree}

The factor \(H_m\) is analytic on \(\Re s>1/2\) for \(m\geq12\), because both dilated zeta
arguments then lie in \(\Re z>6\).  Its nonvanishing is equivalent to
\[
 \zeta(ms)-\zeta((m+1)s)\ne0.
\]
The following quantitative estimate proves more than is needed.

\begin{lemma}[Prime-two domination]\label{lem:two-dominates}
Let \(m\geq12\), \(s=\sigma+it\), and \(\sigma>1/2\).  Put
\[
 A=2^{-ms}-2^{-(m+1)s},
\qquad
 T=\sum_{n\geq3}\bigl(n^{-ms}-n^{-(m+1)s}\bigr).
\]
Then
\[
 |A|>2^{-m\sigma}\frac{29}{100},
 \qquad
 |T|<2^{-m\sigma}\frac{21}{100}.
 \tag{3.1}
\]
\end{lemma}

\begin{proof}
The reverse triangle inequality and \(|2^{-s}|=2^{-\sigma}<2^{-1/2}<71/100\) give
\[
 |A|=2^{-m\sigma}|1-2^{-s}|
 >2^{-m\sigma}\frac{29}{100}.
\]
Set \(x=m\sigma>6\).  For \(n\geq3\),
\[
 |n^{-ms}-n^{-(m+1)s}|
 \leq n^{-x}(1+n^{-\sigma})
 <\frac85 n^{-x}.
\]
Since \(n\geq3\) and \(x>6\),
\[
 n^{-x}=2^{-x}(2/n)^x
 \leq2^{-x}\frac{64}{n^6}.
\]
Consequently
\[
 |T|
 \leq2^{-m\sigma}\frac{512}{5}\sum_{n\geq3}\frac1{n^6}.
\]
The elementary integral estimate
\[
 \sum_{n\geq3}\frac1{n^6}
 \leq\frac1{3^6}+\frac1{4^6}+\frac1{2500}
\]
gives
\[
 \frac{512}{5}
 \left(\frac1{729}+\frac1{4096}+\frac1{2500}\right)
 <\frac{21}{100},
\]
proving the claim.
\end{proof}

\begin{theorem}[Zero-free shell multiplier]\label{thm:zero-free}
For every integer \(m\geq12\), \(H_m\) is holomorphic and zero-free on
\[
 \cD_{1/2}=\{s\in\C:\Re s>1/2\}.
\]
\end{theorem}

\begin{proof}
Absolute convergence gives
\[
 \zeta(ms)-\zeta((m+1)s)=A+T.
\]
Lemma~\ref{lem:two-dominates} implies \(|A|>|T|\), so this difference is nonzero.
Furthermore \(\zeta(ms)\) and \(\zeta((m+1)s)\) are holomorphic and nonzero because their
arguments have real part greater than one.  The identity
\[
 H_m(s)=
 \frac{\zeta(ms)-\zeta((m+1)s)}
 {\zeta(ms)\zeta((m+1)s)}
\]
finishes the proof.
\end{proof}

\begin{corollary}[Exact divisor replication]\label{cor:divisor}
Let \(m\geq12\).  The meromorphic continuation
\[
 S_m^\#(s)=\zeta(s)H_m(s)
\]
has, at every \(s\in\cD_{1/2}\setminus\{1\}\), exactly the same vanishing order as \(\zeta\):
\[
 \ord_s S_m^\#=\ord_s\zeta.
 \tag{3.2}
\]
In particular their zero sets in \(\cD_{1/2}\setminus\{1\}\) agree, with analytic
multiplicities.
\end{corollary}

\begin{proof}
The order of a product of holomorphic functions is the sum of their orders.
Theorem~\ref{thm:zero-free} gives \(\ord_sH_m=0\).
\end{proof}

The integer \(12\) is a uniform threshold delivered by the explicit estimate; no optimality is
claimed.  Its role is to place \(m\Re s\) strictly beyond six throughout \(\cD_{1/2}\).

\section{The \texorpdfstring{\(\Omega\)}{Omega}-leading-echelon law}\label{sec:omega}

We now isolate the general algebraic mechanism.  Let \(\cA\) be the \(\C\)-algebra of arithmetic
functions on the positive integers with Dirichlet convolution
\[
 (f*g)(n)=\sum_{d\mid n}f(d)g(n/d).
\]
It is an integral domain: if \(r\) and \(t\) are the least indices at which nonzero \(f\) and
\(g\) occur, then \((f*g)(rt)=f(r)g(t)\ne0\).

For a prime \(q\), define the prime derivation
\[
 (\partial_qf)(n)=v_q(nq)f(nq).
 \tag{4.1}
\]
A coefficient check gives
\[
 \partial_q(f*g)=(\partial_qf)*g+f*(\partial_qg).
 \tag{4.2}
\]
For \(k\geq0\), write
\[
 L^kf(n)=(-\log n)^k f(n).
 \tag{4.3}
\]
Under Dirichlet transformation, \(L^k\) becomes the \(k\)-th ordinary derivative.

\begin{theorem}[\(\Omega\)-leading-echelon criterion]\label{thm:omega}
Let \(R\) be a finite nonempty set, let \(\ell:R\to\N_{>0}\) be injective, and let
\(f_r\in\cA\).  Suppose that for every \(r\in R\), prime \(q\), and \(k,n\geq0\),
\begin{align}
 \Omega(n)<\ell(r)-1
 &\ \Longrightarrow\ 
 \partial_q(L^kf_r)(n)=0,
 \label{eq:below}
 \tag{4.4}\\
 \Omega(n)=\ell(r)-1
 &\ \Longrightarrow\
 \partial_q(L^kf_r)(n)
 =\ell(r)(-\ell(r)\log q)^k
 \mathbf 1_{\{q^{\ell(r)-1}\}}(n).
 \label{eq:leading}
 \tag{4.5}
\end{align}
Then for every \(J\geq0\), the rectangle
\[
 \{L^kf_r:r\in R,\ 0\leq k\leq J\}
 \label{eq:jet-rectangle}
 \tag{4.6}
\]
is algebraically independent over \(\C\) in \(\cA\).

If, in addition, all \(f_r\) have a finite half-plane of absolute Dirichlet convergence, then
their Dirichlet transforms \(F_r\) are jointly differentially algebraically independent over
\(\C(s)\): the set \(\{F_r^{(k)}:r\in R,k\geq0\}\) is algebraically independent over
\(\C(s)\).
\end{theorem}

\begin{proof}
Fix \(J\), and index the columns by
\[
 I=R\times\{0,\ldots,J\}.
\]
For each \(i=(r,t)\in I\), choose a distinct prime \(q_i\), and form the
\(\cA\)-valued Jacobian
\[
 \mathcal J_{i,(u,k)}
 =\partial_{q_i}(L^kf_u).
\]
Put \(e(u,k)=\ell(u)-1\) and
\[
 N=\prod_{i=(r,t)\in I}q_i^{e(r,t)}.
\]
Every convolution term contributing to \(\det(\mathcal J)(N)\) has total \(\Omega\)-grade
\(\sum_i e(i)\).  By \eqref{eq:below}, no factor can occur below its column grade.  Hence every
factor must lie in its lowest possible grade, and \eqref{eq:leading} computes the coefficient
exactly.

Unique factorization forces a contributing permutation to preserve the equal-\(e\) blocks.
Because \(\ell\) is injective, these are precisely the blocks indexed by one \(u\in R\).
Up to a nonzero scalar, the \(u\)-block is
\[
 \left[(-\ell(u)\log q_{(u,t)})^k\right]_
 {0\leq t,k\leq J}.
\]
Its determinant is a Vandermonde product
\[
 (-\ell(u))^{J(J+1)/2}
 \prod_{0\leq t<v\leq J}
 \bigl(\log q_{(u,v)}-\log q_{(u,t)}\bigr),
\]
which is nonzero because the chosen primes are distinct.  Therefore
\(\det(\mathcal J)(N)\ne0\), so \(\det(\mathcal J)\ne0\) in \(\cA\).

The Jacobian criterion over a characteristic-zero domain now implies algebraic independence:
if a nonzero polynomial relation of least total degree existed, applying the derivations would
place all its nonzero partial derivatives in the kernel, contradicting invertibility of the
Jacobian over \(\Frac(\cA)\).  This proves \eqref{eq:jet-rectangle}.

It remains to justify the coordinate field \(\C(s)\).  The Dirichlet transform is an injective
algebra homomorphism from \((\cA,*)\) into meromorphic germs whenever the sources have finite
abscissa, and it sends \(L^kf\) to the \(k\)-th derivative of the transform.  Suppose, after
clearing denominators, that
\[
 P\bigl(s;(F_r^{(k)}(s))_{r,k}\bigr)=0,
 \qquad P\in\C[s][X_{r,k}]
 \label{eq:coordinate-relation}
 \tag{4.7}
\]
is a nonzero relation involving a finite jet rectangle.  Expand
\(P=\sum_{d=0}^Ds^dP_d\).  Dirichlet multiplication rewrites \eqref{eq:coordinate-relation} as
\[
 \sum_{n\geq1}Q_n(s)n^{-s}=0,
 \qquad
 Q_n(s)=\sum_{d=0}^Ds^d
 \bigl(P_d((L^kf_r)_{r,k})\bigr)(n).
 \tag{4.8}
\]
Uniqueness remains valid for Dirichlet series with uniformly bounded-degree polynomial
coefficients: choose the least \(n\) for which \(Q_n\ne0\), multiply by \(n^s\), and let
\(s=\sigma\to+\infty\); exponential separation dominates the fixed polynomial growth and forces
\(Q_n=0\), a contradiction.  Thus every \(Q_n\) vanishes identically.  Comparing powers of \(s\)
shows \(P_d((L^kf_r)_{r,k})=0\) in \(\cA\) for every \(d\).  The first part forces every
\(P_d=0\), contrary to the choice of \(P\).  Hence no relation over \(\C(s)\) exists.
\end{proof}

\begin{remark}
The proof uses the lowest nonzero associated-graded layer only.  Arbitrary higher
\(\Omega\)-support may be added to each source without changing the conclusion.  This
``upper-tail immunity'' is what makes the criterion applicable to full exponent shells rather
than only to isolated prime powers.
\end{remark}

\section{Zeta and the shells have maximal relative rank}\label{sec:application}

We verify the hypotheses of Theorem~\ref{thm:omega}.  The two elementary facts below are the
entire arithmetic input.

\begin{lemma}[Lowest shell grade]\label{lem:lowest-shell}
If \(a_m(n)\ne0\), then
\[
 \Omega(n)\geq m.
\]
If equality holds, then \(n=p^m\) for a unique prime \(p\).
\end{lemma}

\begin{proof}
Because \(a_m(n)\ne0\), some \(p^m\) divides \(n\), giving \(\Omega(n)\geq m\).
Write \(n=p^mu\).  Additivity of \(\Omega\) gives
\(\Omega(n)=m+\Omega(u)\).  Equality forces \(\Omega(u)=0\), hence \(u=1\).
Uniqueness follows from unique factorization.
\end{proof}

\begin{lemma}[Prime-ray leading layer]\label{lem:prime-ray}
For every \(m\geq1\), prime \(q\), and \(k,n\geq0\),
\begin{align*}
 \Omega(n)<m-1
 &\ \Longrightarrow\
 \partial_q(L^ka_m)(n)=0,\\
 \Omega(n)=m-1
 &\ \Longrightarrow\
 \partial_q(L^ka_m)(n)
 =m(-m\log q)^k\mathbf 1_{\{q^{m-1}\}}(n).
\end{align*}
\end{lemma}

\begin{proof}
If the left side is nonzero, then \(a_m(nq)\ne0\), and
Lemma~\ref{lem:lowest-shell} gives
\[
 m\leq\Omega(nq)=\Omega(n)+1.
\]
This proves the first assertion.  Under equality \(\Omega(n)=m-1\), nonvanishing forces
\(\Omega(nq)=m\), so Lemma~\ref{lem:lowest-shell} gives \(nq=p^m\).  Since \(q\mid p^m\),
one has \(p=q\) and \(n=q^{m-1}\).  At that point
\[
 v_q(q^m)(-\log q^m)^ka_m(q^m)
 =m(-m\log q)^k.
\]
\end{proof}

Let \(\mathbf 1(n)=1\) be the zeta coefficient source.  Assign it level one.  If
\(\Omega(n)=0\), then \(n=1\), and
\[
 \partial_q(L^k\mathbf1)(n)
 =(-\log q)^k\mathbf1_{\{1\}}(n),
\]
which is exactly \eqref{eq:leading} with \(\ell=1\).  There is no lower grade.

\begin{theorem}[Zeta--shell differential freeness]\label{thm:zeta-shell-free}
The countable family
\[
 \{\zeta\}\cup\{S_m:m\geq2\}
\]
is jointly differentially algebraically independent over \(\C(s)\).  Equivalently,
\[
 \{\zeta^{(k)}:k\geq0\}
 \cup
 \{S_m^{(k)}:m\geq2,\ k\geq0\}
 \tag{5.1}
\]
is algebraically independent over \(\C(s)\).
\end{theorem}

\begin{proof}
Take an arbitrary finite subfamily.  Assign level \(1\) to the zeta source \(\mathbf1\) and
level \(m\) to \(a_m\).  The levels are positive and distinct.  The preceding calculation
handles level one, while Lemma~\ref{lem:prime-ray} handles every shell.  All sources are bounded,
so their Dirichlet series converge absolutely on \(\Re s>1\).  Theorem~\ref{thm:omega} applies.
Because every alleged differential-polynomial relation uses only finitely many functions and
derivatives, the finite conclusion proves the countable statement.
\end{proof}

We now separate the common divisor from its transverse directions.  We use the following
standard differential-field observation.

\begin{lemma}[Base adjunction and gauge invariance]\label{lem:gauge}
Let \(K\) be a differential field of characteristic zero.
\begin{enumerate}[label=\textup{(\roman*)}]
\item If \(\{z\}\cup\{y_i:i\in I\}\) is differentially algebraically independent over \(K\),
then \(\{y_i:i\in I\}\) is differentially algebraically independent over
\(K\langle z\rangle\).
\item If \(\{y_i:i\in I\}\) is differentially algebraically independent over \(K\) and
\(a\in K^\times\), then \(\{ay_i:i\in I\}\) is differentially algebraically independent over
\(K\).
\end{enumerate}
\end{lemma}

\begin{proof}
For (i), the set of all derivatives of the \(y_i\) is algebraically independent over the
rational-function field generated by all derivatives of \(z\), by the ordinary tower property
of algebraic independence.

For (ii), for each \(k\),
\[
 (ay_i)^{(k)}
 =a\,y_i^{(k)}
 +\sum_{j<k}\binom{k}{j}a^{(k-j)}y_i^{(j)}.
\]
This is an invertible triangular change of jet coordinates over \(K\), with diagonal entry
\(a\ne0\).
\end{proof}

\begin{theorem}[Relative multiplier hypertranscendence]\label{thm:relative}
Let
\[
 K_\zeta=\C(s)\langle\zeta\rangle.
\]
Then the family
\[
 \{H_m:m\geq2\}
\]
is jointly differentially algebraically independent over \(K_\zeta\).  Equivalently, every
finite subset of
\[
 \{H_m^{(k)}:m\geq2,\ k\geq0\}
 \tag{5.2}
\]
is algebraically independent over \(K_\zeta\).
\end{theorem}

\begin{proof}
Theorem~\ref{thm:zeta-shell-free} and Lemma~\ref{lem:gauge}(i) show that the \(S_m\) are jointly
differentially algebraically independent over \(K_\zeta\).  Since \(\zeta^{-1}\) is a nonzero
element of this meromorphic differential field, Lemma~\ref{lem:gauge}(ii) shows the same for
\(\zeta^{-1}S_m\).  On the initial half-plane of convergence,
\(\zeta^{-1}S_m=H_m\) by Proposition~\ref{prop:factor}; hence the meromorphic germs agree.
\end{proof}

\begin{corollary}[Exact differential rank]\label{cor:rank}
For distinct \(m_1,\ldots,m_r\geq2\) and integers \(J_1,\ldots,J_r\geq0\),
\[
 \trdeg_{K_\zeta}
 K_\zeta\bigl(H_{m_i}^{(k)}:
 1\leq i\leq r,\ 0\leq k\leq J_i\bigr)
 =\sum_{i=1}^r(J_i+1).
 \tag{5.3}
\]
\end{corollary}

\begin{proof}
The displayed generators are algebraically independent by Theorem~\ref{thm:relative}.
\end{proof}

Corollary~\ref{cor:rank} makes ``maximal'' literal: every new shell and every new derivative
raises transcendence degree by exactly one, even after the complete differential data of zeta
have been placed in the coefficient field.

\section{The positive Riemann-divisor frame}\label{sec:rhframe}

Combine Propositions~\ref{prop:partition} and~\ref{prop:factor},
Corollary~\ref{cor:divisor}, and Theorem~\ref{thm:relative}.

\begin{theorem}[Positive hypertranscendental divisor frame]\label{thm:frame}
The family \((S_m)_{m\geq12}\) has all of the following properties.
\begin{enumerate}[label=\textup{(\roman*)}]
\item Each \(S_m\) is a Dirichlet series with coefficients in \(\{0,1\}\), and distinct members
have disjoint coefficient supports.
\item Every coefficient stream is the exact largest-prime-exponent shell
\(\{n:M(n)=m\}\), and all positive shells fuse coefficientwise to
\(\mathbf1-\mathbf1_{\{1\}}\).
\item On \(\Re s>1/2\), each meromorphic continuation has the factorization
\[
 S_m^\#=\zeta H_m
\]
with \(H_m\) holomorphic and zero-free.
\item Away from the pole at \(1\), every \(S_m^\#\) has exactly the zeta zero divisor with
analytic multiplicity.
\item After adjoining \(\zeta\), every derivative of \(\zeta\), and the rational coordinate,
the larger transverse family \((H_m)_{m\geq2}\), and hence its zero-free subfamily
\((H_m)_{m\geq12}\), remains jointly differentially algebraically independent.
\end{enumerate}
\end{theorem}

The theorem exposes two orthogonal structures.  Coefficient supports are disjoint, so no integer
coefficient is reused between two channels.  Analytic zero divisors are identical, so every
right-half-plane zeta zero is reused by all channels.  Relative hypertranscendence says that this
coincidence is exhausted by the common multiplicative factor: it is not enforced by any finite
algebraic differential identity among the zero-free directions.

\begin{corollary}[Exact RH reformulation]\label{cor:rh}
The following statements are equivalent.
\begin{enumerate}[label=\textup{(\roman*)}]
\item The Riemann hypothesis holds.
\item For one \(m\geq12\), \(S_m^\#\) is zero-free in \(1/2<\Re s<1\).
\item For every \(m\geq12\), \(S_m^\#\) is zero-free in \(1/2<\Re s<1\).
\item The common zero set
\[
 \bigcap_{m\geq12}\{s:1/2<\Re s<1,\ S_m^\#(s)=0\}
 \]
is empty.
\end{enumerate}
If RH is false and \(\rho\) is an off-critical zero of multiplicity \(d\), then
\[
 \ord_\rho S_m^\#=d
 \qquad\text{for every }m\geq12.
 \tag{6.1}
\]
\end{corollary}

\begin{proof}
Corollary~\ref{cor:divisor} identifies the zero set of every \(S_m^\#\) in the strict right
half-plane with that of \(\zeta\), including multiplicity.  The functional equation and Schwarz
symmetry make absence of zeros in \(1/2<\Re s<1\) equivalent to RH; see, for example,
Titchmarsh~\cite[Chapter II]{Titchmarsh1986}.
\end{proof}

\subsection{What the frame adds to the RH problem}

Corollary~\ref{cor:rh} alone would be only an equivalent reformulation.  The content of
Theorem~\ref{thm:frame} is the additional structure attached to it.

First, false RH would force the same analytic cancellation independently inside every set
\[
 \{n:M(n)=12\},\ \{n:M(n)=13\},\ \ldots,
 \tag{6.2}
\]
although these sets are pairwise disjoint and their indicators are nonnegative.  This turns one
global zero into a countable simultaneous-cancellation constraint.

Second, normalization does not collapse these constraints into a finite differential system.
Even over the complete zeta differential field, any finite packet of normalized channels and
any finite collection of their derivatives has full algebraic rank.  A proof based on
simultaneous channels must therefore use analytic information---positivity, size, phase,
boundary behavior, or interaction across infinitely many levels---rather than a hidden finite
differential identity.

Third, the \(\Omega\)-leading-echelon criterion is not specific to zeta zeros.  It applies to any
Dirichlet-source family whose selected-prime derivatives have distinct lowest
total-prime-factor grades and pure prime-ray leading terms.  Thus the same mechanism can
separate a prescribed common divisor from a maximally free transverse Dirichlet frame in other
Eulerian settings.

These statements answer the structural question posed in the introduction.  Exact positive
fusion, complete common-divisor replication, and maximal relative differential rank are
compatible, and the exponent-shell family realizes all three unconditionally.

\section{Conclusion}

The largest-exponent partition of the integers produces a rigid analytic object that is not
visible from its elementary definition.  Every high shell is a positive Dirichlet subseries of
the zeta coefficient stream.  Every shell from level twelve onward reproduces the entire
strict-right-half zeta zero divisor.  After the common zeta factor is removed, the larger
multiplier family beginning already at level two is countably differentially independent over
the full zeta differential field.

The universal reason is the \(\Omega\)-leading layer.  Selected-prime derivations turn distinct
exponent levels into distinct pure-prime grades; finite jet Jacobians reduce to nonzero
prime-log Vandermonde blocks, and all remaining arithmetic support lies too high to cancel
them.  This yields an exact common-factor/transverse-freeness law:
\[
\boxed{\begin{gathered}
 \text{one Riemann divisor}+\text{disjoint positive sources}\\
 {}+\text{countably maximal transverse differential rank}
\end{gathered}}
\]
For RH, the outcome is a new positive divisor frame with a precise failure signature:
every off-critical zero must recur with unchanged multiplicity in every disjoint shell.  The
paper proves that structure completely; it neither assumes nor asserts the absence of such
zeros.

\begin{thebibliography}{99}

\bibitem{GrosswaldSchnitzer1978}
E.~Grosswald and F.~J.~Schnitzer,
\newblock A class of modified \(\zeta\) and \(L\)-functions,
\newblock \emph{Pacific J. Math.} \textbf{74} (1978), 357--364.
\newblock \href{https://doi.org/10.2140/pjm.1978.74.357}
{doi:10.2140/pjm.1978.74.357}.

\bibitem{Kolchin1973}
E.~R.~Kolchin,
\newblock \emph{Differential Algebra and Algebraic Groups},
\newblock Academic Press, New York, 1973.

\bibitem{Nagoshi2016}
H.~Nagoshi,
\newblock The hypertranscendence of \(L\)-functions for
\(\mathrm{GL}_m(\mathbb A_{\mathbb Q})\),
\newblock \emph{Bull. Aust. Math. Soc.} \textbf{93} (2016), 388--399.
\newblock \href{https://doi.org/10.1017/S000497271500129X}
{doi:10.1017/S000497271500129X}.

\bibitem{Ostrowski1920}
A.~Ostrowski,
\newblock \"Uber Dirichletsche Reihen und algebraische Differentialgleichungen,
\newblock \emph{Math. Z.} \textbf{8} (1920), 241--298.
\newblock \href{https://doi.org/10.1007/BF01206523}
{doi:10.1007/BF01206523}.

\bibitem{Pong2016}
W.~Y.~Pong,
\newblock Applications of differential algebra to algebraic independence of arithmetic
functions,
\newblock \emph{Acta Arith.} \textbf{172} (2016), 149--173.
\newblock \href{https://doi.org/10.4064/aa8112-12-2015}
{doi:10.4064/aa8112-12-2015}.

\bibitem{Postnikov1949}
A.~G.~Postnikov,
\newblock On the differential independence of Dirichlet series,
\newblock \emph{Dokl. Akad. Nauk SSSR} \textbf{66} (1949), 561--564
\newblock (in Russian).

\bibitem{ShapiroSparer1986}
H.~N.~Shapiro and G.~H.~Sparer,
\newblock On algebraic independence of Dirichlet series,
\newblock \emph{Comm. Pure Appl. Math.} \textbf{39} (1986), 695--745.
\newblock \href{https://doi.org/10.1002/cpa.3160390602}
{doi:10.1002/cpa.3160390602}.

\bibitem{Titchmarsh1986}
E.~C.~Titchmarsh,
\newblock \emph{The Theory of the Riemann Zeta-Function}, 2nd ed.,
revised by D.~R.~Heath-Brown,
\newblock Oxford University Press, Oxford, 1986.

\end{thebibliography}

\end{document}
